\documentclass[tikz,border=6pt]{standalone}
\usepackage{peabra}
\begin{document}
\begin{tikzpicture}[x=1cm,y=1cm, a/.style={doboz=24mm, font=\footnotesize, minimum height=15mm}, b/.style={sotetdoboz=24mm, font=\footnotesize, minimum height=15mm}]
  \foreach \i/\t/\s in {1/{1. tervezés és\\előkészítés}/a, 2/{2. struktúra-\\elemzés}/a, 3/{3. funkció-\\elemzés}/a, 4/{4. hiba-\\elemzés}/a, 5/{5. kockázat-\\elemzés}/b, 6/{6. optimalizálás}/b, 7/{7. eredmények\\dokumentálása}/a} {
    \node[\s] (n\i) at ({(\i-1)*2.85},0) {\t};
  }
  \foreach \a/\b in {1/2,2/3,3/4,4/5,5/6,6/7} \draw[nyil] (n\a) -- (n\b);
  \draw[peKek, line width=0.6pt, decorate, decoration={brace, amplitude=5pt, mirror}] (-1.2,-1.0) -- (1.2,-1.0) node[midway, below=6pt, felirat] {hatókör, csapat, ütemezés};
  \draw[peKek, line width=0.6pt, decorate, decoration={brace, amplitude=5pt, mirror}] (1.65,-1.0) -- (9.75,-1.0) node[midway, below=6pt, felirat, align=center] {rendszerelemzés: elemek, funkciók, hibaláncok\\(hibahatás -- hibamód -- hibaok)};
  \draw[peKiemel, line width=0.6pt, decorate, decoration={brace, amplitude=5pt, mirror}] (10.2,-1.0) -- (15.6,-1.0) node[midway, below=6pt, felirat, align=center, text=peKiemel] {hibaelemzés és kockázatcsökkentés:\\S, O, D, prioritás, intézkedések};
  \draw[nyil, draw=peKiemel, dash pattern=on 3pt off 2pt] (n6.north) -- ++(0,0.75) -| node[felirat, pos=0.25, above, text=peKiemel] {új értékelés az intézkedések után} (n5.north);
\end{tikzpicture}
\end{document}
